\documentclass[conference,9pt]{IEEEtran}

% -------------------- Packages --------------------
\usepackage{graphicx}
\usepackage{amsmath, amssymb}
\usepackage{booktabs}
\usepackage{multirow}
\usepackage{array}
\usepackage{caption}
\usepackage{subcaption}
\usepackage{float}
\usepackage{url}
\usepackage[hidelinks]{hyperref}

% Optional: better table column widths
\newcolumntype{L}[1]{>{\raggedright\arraybackslash}p{#1}}

% -------------------- Title + Authors --------------------
\title{\textbf{ResNet18: An Interpretable Deep Learning Framework for Automated Eye-Disease Classification}

\author{
\begin{tabular}{ccc}
\textbf{Afrina Alam Lamia} &
\textbf{Ohiduzzaman Rayhan} &
\textbf{Sumiaya Sharif} \\

ID: 2022-1-60-231 &
ID: 2023-1-60-108 &
ID: 2023-1-60-048 \\

Department of CSE &
Department of CSE &
Department of CSE \\

East West University &
East West University &
East West University \\

\texttt{afrinaaa0711@gmail.com} &
\texttt{ohiduzzamanrayhan866@gmail.com} &
\texttt{sumiayasharif256@gmail.com}
\end{tabular}
}


}

\begin{document}
\maketitle

% -------------------- Abstract --------------------
\begin{abstract}
Early and accurate eye disease detection is crucial, as observational methods are time-consuming and susceptible to human error. To mitigate these issues, we present a deep-learning based robust framework for multi-class eye disease classification using RGB images. Identified ResNet18 as the optimal model, after evaluating multiple architectures including CNNs, vision transformers and CNN-transformer hybrids. The ResNet18 model achieved a test accuracy of 97.44\% and outperformed other models. The proposed pipeline incorporates image pre-processing (grayscale conversion, CLAHE, Gaussian filtering), data augmentation, and class-balancing using SMOTE. Grad-CAM visualization confirmed that the model focuses on clinically relevant ocular regions, improving interpretability. Extensive evaluation using cross-validation, confidence intervals, and statistical tests verified the reliability and robustness of the approach. This framework shows strong potential as an automated diagnostic tool for eye disease classification within its defined classes.
\end{abstract}

% Optional keywords
\begin{IEEEkeywords}
Medical image classification, Deep learning, CNN, Vision Transformer, Grad-CAM, Eye disease detection
\end{IEEEkeywords}

% -----
% -------------------- 2. Literature Review --------------------
\section{\textbf{Literature Review}}



D.M. et al. \cite{1} in this work used four pre-trained CNN models to achieve accurate multiclass classification of eye diseases from external eye images. Four pre-trained models were VGG16, ResNet50, InceptionV3 and MobileNetV2. These models were fine tuned using ImageNet weights. Adam Optimizer, cross-entropy loss were used to train the models and 10 epochs for ReLU activation. The confusion metrices and evaluation metrices (F1- score, recall, precision, accuracy) were used to evaluate and compare the performance of the models. Best accuracy was achieved by ResNet50. \\

M.A. et al. \cite{2} have proposed a Multi-Level Deep Convolutional Neural Network (ML-DCNN) for glaucoma detection and categorisation based on retinal images. The whole network comprises two different CNN designs, namely Detection-Net CNN for glaucoma detection and Classification-Net CNN for glaucoma categorisation into Normal, Early, Moderate, and Advanced glaucoma. The network is optimised using an Adam optimiser with a cross-entropy loss function for model evaluation using accuracy, sensitivity, specificity, and precision, along with a confusion matrix. The proposed ML-DCNN demonstrated a high accuracy of 99.39\%, which is better than other existing solutions.\\

Zhou et al. \cite{3} proposed a ViT-based model with multiclass classification of retinal diseases using OCT images from the Kermany2018 Kaggle dataset, with NORMAL, CNV, DME, and DRUSEN classes. Features were learned in this model using image patching, token embedding involving positional encoding, and multi-head self-attention. In developing the model, a symmetric cross-entropy loss function was applied to reduce overfitting noise. It was evaluated on the basis of accuracy, sensitivity, specificity, and confusion matrices with 96.57\% validation accuracy and 95.76\% test accuracy over EfficientNet, ResNet50, and VGG16 models.\\

Akça et al.\cite{4} performed a comparative analysis of vision transformer architectures for multiclass classification of retinal diseases from OCT images using transformers. A balanced subset of the Labelled OCT dataset with CNV, DME, Drusen, and Normal classes was utilised. Among the attention-driven models, Vision Transformer (ViT), T2T-ViT, and Mobile-ViT were tested using the AdamW optimizer with cross-entropy loss and standard evaluation metrics such as accuracy, precision, recall, F1-score, and confusion matrices. Among all models, the highest accuracy of 99.17\% was obtained using Mobile-ViT.\\

Mondal et al.\cite{7} in this study employed three pre-trained CNNs for the identification of conjunctivitis in external eye images. There were more pre-trained models used in the research, including InceptionV3, ResNet-50, and VGG-19. The models were fine-tuned using weights from ImageNet. The models were trained using the Adam and Adadelta optimisers and cross-entropy loss functions. Hyperparameter search using Keras Tuner was used for model optimisation. DBSCAN algorithm clustering was utilised for grouping training images with homogeneous features into clusters. Each distinct group required a unique model. The models were validated using metrics of accuracy, recall rate, F1 score, entropy, and AUC-ROC. InceptionV3 resulted in a higher accuracy score compared with the other models. ResNet-50 and VGG-19 models resulted in favourable results.\\

Nakayama et al. \cite{8} introduced CAD-EYE, a hybrid deep learning model for the classification of eye diseases, including diabetic retinopathy, hypertensive retinopathy, glaucoma, and contrast-related problems. The CAD-EYE model is a combination of MobileNetV2 and EfficientNetB0 with feature fusion and uses an XGBoost classifier. The model used global average pooling and dense layers with ReLU and SoftMax activation functions and trained using the Adam optimiser and cross-entropy function for 10 epochs. Accuracy, F1-score, recall, and precision were used for evaluation, and CAD-EYE reported 98\% accuracy, which outperformed other models, including VGG16, ResNet50, and InceptionV3.\\

A brief literature survey on deep learning, especially on CNNs, is given by Alzubaidi et al. \cite{9}. Topics include basics and notable CNN architectures (LeNet, AlexNet, VGG, GoogLeNet, ResNet, DenseNet, and EfficientNet) and discuss challenges and future work without training any model and without any experiments (no data, optimiser, loss function, epochs, or accuracy given), only a literature survey across computer vision, medical imaging, speech, NLP, and autonomous vehicles with a qualitative comparison by previous works. Challenges include requirements for data, overfitting, computational efficiency, interpretability, vanishing and exploding gradients, and applicability in low-resource environments. Ultimately, their findings support that future research on CNN could be directed towards efficiency, explainability, Hybrid models, and less dependence on a large amount of data for training and annotation purposes.\\

Babaqi et al.\cite{13} applied CNNs with transfer learning to classify eye diseases using colour images. They created a set of approximately 4200 images categorised by eye conditions into labels Normal, Cataract, Diabetic Retinopathy, and Glaucoma. Results were measured using accuracy, precision, recall, F1, and accuracy. Results showed a higher accuracy of 94\% using transfer learning, contrasted with 84\% accuracy in a traditional CNN. This paper indicates that transfer learning improves eye disease image classification accuracy, with difficulties being few images in larger datasets, along with a need to use more types of eye diseases.\\

Ayesha et al. \cite{14} proposed a CNN model integrated with a Vision Transformer-based Self-Attention Mechanism (SAM) for the classification of eye diseases from pictures of the external eye. Based upon 4,217 images and for four classes, the CNN-ViT model had five convolutions involving ReLU activation and max pooling layers, along with a 16-head SAM and 64 key dimensions and three dense layers involving softmax activation. It achieved 94\% accuracy and an average AUC of 98.82\% upon Adam optimisation and sparse categorical cross-entropy in 25 epochs and was evaluated through the use of confusion matrices and metrics (accuracy, precision, recall, F1 score, and AUC ROC). The result for Diabetic Retinopathy was perfect for precision, recall, and F1 score.\\

Topaloglu et al.\cite{15} developed a CNN architecture with a unique Care Model for Multiclass eye diseases.Classification of eye diseases using retinal images. With a Kaggle competition database of 45,231 high-resolution images of diabetic retinopathy. It consists of convolution layers, max pooling, two dense layers, and a softmax function. It uses a Care Model, which produces new learnable images by using pixel-wise filtering and Global Average Pooling. Optimization was done using Adam and categorical cross-entropy loss functions in 50 epochs. Validation is done via confusion matrices with accuracy, precision, recall, and F1. It produced a test accuracy of 87\% accuracy.\\

Jain et al. \cite{16} proposed the Adaptive Hybrid Net-CNN: a binary \& multi-class eye disease detection using a deep hybrid model with InceptionV3, MobileNetV3-Small, ResNet-101, \& Adaptive CNN layers. The binary dataset split is train/test = 80/20 while multi-class datasets split into three classes \& SMOTE. There were a total of 50 epochs with a batch size of 132. The evaluation metrics used were acc, F1-score, precision, recall, sensitivity \& specificity. The model resulted in a higher accuracy rate \& sensitivity with 98.45\% \& 96.36\%, respectively. However, more testing is needed beyond APTOS-2019.\\

Lin et al. \cite{17}employed a Two-StreamCNN-Transformer architecture for the classification of retinal disease. Spatial features were modeled by CNNs, and transformers interpreted overall information. The Adam optimiser and cross-entropy loss function were applied with a batch size of 16 and a total of 50 epochs, and data augmentation techniques of rotation, flip, cropping, and adding Gaussian noise were employed. Accuracy of 96\%, F1-score of 0.96, and an AUC of 1.00 were achieved by LGSF-Net, outperforming ResNet50, ViT, Swin, and InceptionV3 models while utilising lower parameters of 18.7K and FLOPs of 0.93G.\\

% -------------------- 3. Methodology --------------------
\section{\textbf{Methodology}}
\vspace{-0.3em}
In this section, we have discussed the methods used for the classification of eye-disease using simple RGB eye images. The overall approach is demonstrated in Fig. 1, which illustrates the workflow from data collection to correct prediction of eye diseases among five distinctive classes.
\begin{figure}[H]
    \centering
    \includegraphics[width=\columnwidth]{figures/workfloww.jpeg}
    \caption{Overall workflow of the proposed eye disease classification framework.}
    \label{fig:workflow}
\end{figure}


\subsection{Dataset Description }
    For eye-disease classification study, we have used a publicly available dataset from Mendeley Data repository \cite{dataset} This particular dataset is designed to automate eye disease detection using image-based analysis. It consists of real-world clinical RGB eye-images of five categories: Cataract, Conjunctivitis, Eyelid, Normal and Uveitis – each containing images after SMOTE (Synthetic Minority Over-sampling Technique) was applied, ensuring a balanced distribution among classes. 
    
\begin{table}[H]
\centering
\caption{Class-wise Data Distribution.}
\label{tab:dataset_dist}
\begin{tabular}{l c}
\toprule
\textbf{Class} & \textbf{Image Count} \\
\midrule
Cataract & 544 \\
Conjunctivitis & 357 \\
Eyelid & 525 \\
Normal & 649 \\
Uveitis & 223 \\
\bottomrule
\end{tabular}
\end{table}
% Example figure slot
\begin{figure}[H]
\centering
\includegraphics[width=\linewidth]{figures/class_sampless.png}
\caption{Sample images per class from the dataset.}
\label{fig:class_samples}
\end{figure}

\subsection{Data Preprocessing}
     A preprocessing pipeline was applied to all raw images for image quality enhancement and noise suppression. The preprocessing steps standardize image resolution by resizing, along with enhancing local contrast and smoothing filtering to ensure the output is compatible with deep learning models. 
In the process, OpenCV was used to read images for invalid image handling. Since OpenCV uses the BGR channel order, BGR to RGB conversion was performed for overall consistency in operation. For stable training and efficient computation, all images were resized to 224×224 pixels, a fixed uniform input resolution. After resizing, the RGB images were transformed to grayscale to reduce color dependency. This step brings attention to intensity-based structural representation. To boost local contrast and visibility of minor features, Contrast Limited Adaptive Histogram Equalization (CLAHE) \cite{clahe} was implemented on the grayscale images, configured with a clip threshold of 2.0 and a grid partition of 8x8 tiles. This limitation strategy reduces noise exaggeration while improving visibility in low-contrast areas. Following contrast enhancement, a Gaussian filtering using 3x3 convolution kernel was employed to minimize high-frequency noises. 
As a final step, two of widely known image quality metrices – Peak Signal to Noise Ratio (PSNR) and Structural Similarity Index (SSIM) were computed on first 200 processed images. It attained an average PSNR of 21.40 and SSIM of 0.853, indicating minimal structural distortion. 
Fig.3 previews the paired samples of before and after preprocessing. The outcomes confirms the structural preservation and medical relevance without excessive distortion.


% Figure: Raw vs Preprocessed
\begin{figure}[H]
\centering
\includegraphics[width=\linewidth]{figures/raw_vs_preprocessed.png}
\caption{Raw vs. preprocessed image samples}
\label{fig:raw_pre}
\end{figure}

\subsection{Dataset Splitting and Augmentation}
Post-preprocessing stage, the dataset was segmented into train, validation and test subsets to support stable model learning and fair model evaluation. A fresh split directory was created to prevent data leakage from earlier data. Images were randomly shuffled and partitioned using 80:10:10 training, validation and testing ratio, applied independently on each class (Fig.4). The resulting sets were then organized into a class-specific folder hierarchy, which supports PyTorch-compatible dataset loading. 

\begin{table}[H]
\centering
\caption{Train/Validation/Test split per class.}
\label{tab:split}
\begin{tabular}{l c c c}
\toprule
\textbf{Class} & \textbf{Train} & \textbf{Val} & \textbf{Test} \\
\midrule
Cataract & 435 & 54 & 55 \\
Conjunctivitis & 285 & 35 & 37 \\
Eyelid & 420 & 52 & 53 \\
Normal & 519 & 64 & 66 \\
Uveitis & 178 & 22 & 23 \\
\bottomrule
\end{tabular}
\end{table}

\begin{figure}[H]
    \centering
    \includegraphics[width=\columnwidth]{figures/data_split.png}
    \caption{Class-wise distribution of images across training, validation, and test splits.}
    \label{fig:data_split}
\end{figure}
Data augmentation was exclusively performed on the training set only to promote generalization while preventing overfitting effects. Inputs were uniformly rescaled to 224x224 and subjected to stochastic transformations by random horizontal flipping and minor angular rotations (±10°). In contrast, the validation and test sets left unchanged to ensure fair and unbiased evaluation. Finally, the entire dataset was transformed into tensors followed by normalization using ImageNet statistical values, enabling consistent input scaling for learning pretrained models. We designed a custom-made PyTorch Dataset that retrieves images from the split folders and maps class names into integer labels. DataLoaders were configured with mini-batches of size of 32 to optimize batch-wise training. The shuffling was applied only to the training loaders when validation and test loaders remained unshuffled. 
In summary, the splitting and augmentation approach facilitates unbiased performance assessments, prevents overfitting issues through training-only augmentation and supports and efficient deep-learning workflow.


\subsection{Model Architectures and Training Strategy}


\subsubsection{Evaluated CNN and Transformer Models}

To systematically evaluate the efficacy of CNN \cite{resnet} and Vision Transformer architectures \cite{vit,swin,mobilevit,levit}  for clinical eye disease classification, a range of pretrained models were trained with standardize preprocessing and training settings. The models are as follows:

\begin{itemize}
    \item \textbf{Swin-Tiny}:A hierarchical vision transformer model that computes self-attention inside overlapping local windows supporting both global and local image features. An ImageNet-pretrained Swin-Tiny network was fine-tuned for eye-disorder diagnosis by replacing the output head with a custom classification layer. Dropout regularization was also employed to limit overfitting.

    \item \textbf{DeiT-Tiny}: A data-efficient Image transformer suited for scares data scenerios through efficient training mechanisms. In our process, the DeiT-Tiny architecture initialized with pretrained weights was fine-tuned on the processed data to evaluate the feasibility of small-scale transformer networks.

    \item \textbf{LeViT-256}: A CNN-Transformer hybrid architecture fusing convolutional feature biases with lightweight self-attention layers to reduce computational complexity and minimize inference cost. The pretrained LeViT-256 backbone preserved its high representational capacity by reconfiguring the final fully connected layer and trained with early stopping to improve generalization in medical image classification.

    \item \textbf{MobileViT-XS}: MobileViT-XS fuses lightweight convolutional layers with attention-based transformer modules to learn both local and global context while maintaining computational efficiency. This model was included to examine transformer behavior in low-resource scenarios and was fine-tuned using ophthalmic dataset, following the same preprocessing and training pipeline. 

    \item \textbf{ShuffleNetV2}: A compact parameter-efficient convolutional neural network (CNN) by channel shuffling operations alongside pointwise grouped convolution layers. It functioned as a baseline CNN model with its output head reconfigured by a dropout layer followed by a linear layer to handle five distinctive disease categories.
\end{itemize}

\subsubsection{\textbf{Proposed ResNet18 Model}}

Among all evaluated architectures, ResNet18 consistently outperformed other models’ performances reflecting robust generalization under all reported metrices. 
The model follows residual learning framework, where identity shortcut connections enable stable gradient propagation across deep network layers. This characteristic is highly suited for medical image analysis, where subtle textural patterns and detailed visual cues must be retained. In this study, the ImageNet-pretrained ResNet18 was modified by replacing the final FC layer with a regularized linear classification head to limit model overfitting. Label smoothing was also introduced within cross-entropy loss formulation to reduce over-confident predictions and applied AdamW optimization with weight decay to promote model generalization. Later on, gradient magnitude clipping and adaptive learning rate scheduling were employed to ensure stable training. The combination of robust CNN inductive biases, strong regularization techniques and transfer-based optimization outperform both CNNs and transformer model architectures on the eye-disease prediction task. 
Fig. 5 illustrates the architecture of ResNet18-CNN Model \cite{resnet}
\begin{figure}[H]
    \centering
    \includegraphics[width=\columnwidth]{figures/resnet_arch1.drawio.png}
    \caption{ResNet18- a CNN-based Model Architecture}
    \label{fig:resnet_arch}
\end{figure}

\subsection{Data Training Strategy}
       All experiments were conducted using PyTorch deep learning framework performed on GPU-enabled hardware to guarantee efficient and reproducible training. Input images were resized to $224 \times 224$ resolution normalized following ImageNet conventions to satisfy pretrained network input requirements. The models were trained using AdamW optimizer with a base learning rate of $3 \times 10^{-4}$ and a weight decay coefficient of $1 \times 10^{-4}$ to improve generalization. Categorical cross-entropy loss incorporating label smoothing of 0.1 was employed to prevent excessive confidence and enhance model robustness under class ambiguity. A ReduceLROnPlateau scheduler was applied based on validation loss, progressively reducing the learning rate when performance plateaued. Some regularization strategies were applied including dropout (rate=0.5), gradient clipping (max norm = 5.0) and training-time data augmentation. A mini-batch size of 32 was used throughout with random shuffling,flipping and rotation. All networks trained with a maximum epoch limit of 40 using early stooping with a patience value of 5. The best performing validation model was chosen for final testing on the test split. This standardized training strategy confirms unbiased and controlled comparison while maximizing task performance.
\begin{table}[H]
\centering
\caption{Training configuration and hyperparameter settings for the proposed model.}
\label{tab:training_config}
\begin{tabular}{l l}
\toprule
\textbf{Parameter} & \textbf{Value} \\
\midrule
Model & ResNet18 (ImageNet pretrained) \\
Input Size & $224 \times 224 \times 3$ \\
Number of Classes & 5 \\
Optimizer & AdamW \\
Learning Rate & $3 \times 10^{-4}$ \\
Weight Decay (L2) & $1 \times 10^{-4}$ \\
Loss Function & CrossEntropyLoss \\
Label Smoothing & 0.1 \\
Batch Size & 32 \\
Maximum Epochs & 40 \\
Early Stopping & Yes (patience = 5) \\
Learning Rate Scheduler & ReduceLROnPlateau \\
Scheduler Factor & 0.5 \\
Scheduler Patience & 2 \\
Gradient Clipping & $\text{max\_norm} = 5.0$ \\
Dropout Rate & 0.5 \\
Framework & PyTorch \\
Device & GPU \\
\bottomrule
\end{tabular}
\end{table}


\subsection{Evaluation Metrics and Statistical Tests }
       Model performance was measured using both traditional and domain-specific metrices. In addition to total accuracy, class-wise evaluation metrices including Sensitivity, Specificity, Positive Predictive Value (PPV), and Negative Predictive Value (NPV) were derived to examine diagnosis performance. For classification accuracy, bootstrap-based Confidence Interval (CI) was computer using iterative resampling of the test set. For statistical comparison, McNemar statistical test was performed to evaluate paired prediction \cite{mcnemar}. It confirms whether performance variations are significantly meaningful. These assessment strategies deliver a thorough evaluation of models’ classification effectiveness.

\subsection{Explainable AI using Grad-CAM}
For better understanding of the proposed models, Gradient Weighted Class Activation mapping (Grad-CAM) was generated for top-performing models including ResNet18 and MobileViT-XS. Grad-CAM produces category-focused heatmaps by analyzing gradients into the final feature extraction layers, highlighting the most relevant regions for classification results. This process allows visual inspection of whether the networks focus on medically relevant eye regions instead or irrelevant non-informative background. This exploratory analysis improves confidence in predictions and enhances model credibility and transparency.

% -------------------- 4. Results and Discussion --------------------
\section{\textbf{Results and Discussion}}

\subsection{Model Performance Comparison}
% Add your performance summary table
\begin{table}[H]
\centering
\caption{Performance comparison of different models on training, validation, and test sets.}
\label{tab:model_performance}
\begin{tabular}{l c c c}
\toprule
\textbf{Model} & \textbf{Train Accuracy} & \textbf{Validation Accuracy} & \textbf{Test Accuracy} \\
\midrule
ResNet18      & 1.0000 & 0.9604 & 0.9744 \\
Swin-Tiny     & 0.9951 & 0.9559 & 0.9444 \\
DeiT-Tiny     & 0.9858 & 0.9604 & 0.9359 \\
LeViT-256     & 0.9989 & 0.9780 & 0.9402 \\
ShuffleNetV2  & 0.9989 & 0.9648 & 0.9359 \\
MobileViT-XS  & 0.9940 & 0.9648 & 0.9573 \\
\bottomrule
\end{tabular}
\end{table}
Table IV reports performance across all data splits (training, validation, test set) of all considered architectures. Among the six evaluated networks, ResNet18 obtained the highest testing accuracy of 97.44\%, indicating stronger generalization ability. Although the model has near-saturated training results, the minor performance gap between training, validation (96.04), and test accuracy confirms that overfitting was well controlled by means of explicit regularization and early stopping. LeVit-256 attained high validation accuracy (97.80\%) but showed a noticeable drop in test performance pointing to moderate overfitting. ShuffleNetV2 and DeiT-Tiny performed well during training but had lower generalization on test data (93.59\%) indicating that compact architectures may have difficulty capturing subtle ocular features with restricted dataset sizes. Vision transformer-based models such as Swin-Tiny and DeiT-tiny showed strong training performance but reflected relatively lower accuracy during evaluation and testing, highlighting shortcomings in extracting class-specific patterns from existing dataset scale. The MobileViT-XS hybrid model outperformed transformer-only models, achieving a 95.73\% accuracy on test set, demonstrating  the benefits of incorporating convolutional feature biases with transformer attention layers. Overall analysis shows CNN architectures delivered more stable and consistent accuracy across all data splits, making ResNet18 the most robust model for ocular disease classification.
\begin{figure}[H]
    \centering
    \includegraphics[width=\columnwidth]{figures/comparison.png}
    \caption{Model-wise train,validation and test accuracy plot. }
    \label{fig:comparison}
\end{figure}

\subsection{Accuracy and Loss Curve Analysis}
% Two figures side by side
\begin{figure}[H]
\centering
\begin{subfigure}{0.48\linewidth}
\centering
\includegraphics[width=\linewidth]{figures/acc_curve1.png}
\caption{Accuracy curve.}
\end{subfigure}
\hfill
\begin{subfigure}{0.48\linewidth}
\centering
\includegraphics[width=\linewidth]{figures/loss_curve.png}
\caption{Loss curve.}
\end{subfigure}
\caption{Training and validation learning curves}
\label{fig:curves}
\end{figure}
The training and validation accuracy curve against epochs of the best performing model (ResNet18) exhibit early learning convergence during early epochs, with training accuracy crossing 95\% in the initial training iterations because of ImageNet pretrained weights. As training progress, the model attains almost perfect training accuracy while preserving stable validation performance around 96-97\%. The tight alignment between training and validation curve underscores the effectiveness of stabilization techniques. 
Similarly, the loss curve exhibits a steady decrease with training loss stabilizing around 0.40 and validation loss near 0.46. No instability or sudden increase in validation loss is noticed, confirming convergence stability and controlled learning behavior. The ultimate test accuracy of 97.44\% closely aligns with validation accuracy, supporting the robustness of the optimized model on unseen data.

\subsection{Misclassification Analysis}
\begin{figure}[H]
    \centering
    \includegraphics[width=\columnwidth]{figures/cmm.png}
    \caption{Confusion Matrix for best performing model.}
    \label{fig:cmm}
\end{figure}
Fig. 8 illustrates the confusion matrix visualization for the best performing model - ResNet18 analyzed on the test dataset. It shows highly accurate prediction for cataract, eyelid, and Normal classes with zero misclassification cases, reflecting strong class separation. The majority of errors for ResNet18 are concentrated in the Uveitis samples, where a few instances are predicted incorrectly. 3 samples of Uveitis are wrongly predicted as Conjunctivitis which shows a minor contradiction between Conjunctivitis and Uveitis classes. It might happen in case of visual resemblance among inflammatory ocular conditions. However, the prediction model maintain strong diagonal dominance confirming reliable multi-class classification, while identifying Uveitis as the most error-prone class due to inter-class visual overlap. 

\subsection{Evaluation Metrics Comparison }
\begin{table}[H]
\centering
\caption{Comparison of test accuracy, confidence intervals, and macro-averaged medical evaluation metrics across models.}
\label{tab:advanced_metrics}
\begin{tabular}{l c c c c c}
\toprule
\textbf{Model} &
\textbf{CI$_{95\%}^{Low}$} &
\textbf{CI$_{95\%}^{High}$} &
\textbf{Sensitivity} &
\textbf{Specificity} &
\textbf{PPV} \\
\midrule
ResNet18 & 0.9530 & 0.9915 & 0.9511 & 0.9936 & 0.9738 \\
MobileViT-XS & 0.9316 & 0.9829 & 0.9378 & 0.9898 & 0.9379 \\
Swin-Tiny & 0.9145 & 0.9744 & 0.9249 & 0.9869 & 0.9202 \\
LeViT-256 & 0.9103 & 0.9701 & 0.9212 & 0.9852 & 0.9296 \\
DeiT-Tiny & 0.9060 & 0.9658 & 0.9347 & 0.9851 & 0.9135 \\
ShuffleNetV2 & 0.9060 & 0.9658 & 0.9211 & 0.9847 & 0.9163 \\
\bottomrule
\end{tabular}
\end{table}
Table V reports detailed comparison of all trained architectures using diagnostic evaluation metrices, including 95\% confidence intervals (CI), macro-averaged sensitivity, specificity, PPV and NPV. Among all, ResNet18 outperformed all other models on the test split (97.44\%), with the tightest CI (95.30\%-99.15\%), indicating robust generalization and statistical stability. ResNet18 also maintained the strongest balance between sensitivity (95.11\%) and specificity (99.36\%), which is crucial for limiting both false-negative and false-positive errors in clinical diagnosis.
MobileViT-XS achieved second-best performance, particularly in specificity (98.98\%) and NPV (98.98\%), confirming its robustness despite its parameter-efficient design.  Vision transformer models such as Swin-Tiny, LeViT-256, and DeiT-Tiny delivered competitive but slightly lower results , likely due to their strong data dependency. While, ShuffleNetV2 showed efficient computation but modestly decreased sensitivity and PPV, indicating increased misclassification risk in less frequent disease categories. 


% Optional McNemar results figure
\begin{figure}[H]
\centering
\includegraphics[width=\linewidth]{figures/mcnemar.png}
\caption{McNemar paired test to evaluate the statistical significance of the best model.}
\label{fig:mc}
\end{figure}
To systematically assess whether the performance gaps were statistically justified, the McNemar paired test was performed using test-set outputs. The comparison between ResNet18 and MobileViT-XS produced a p-value of 0.289, indicating equivalent statistical performance, even though ResNet18 attained overall higher accuracy. In comparison, the test between ResNet18 and ShuffleNetV2 returned a statistically significant p-value of 0.022, validating a significant performance improvement in the best model. These findings clearly justify ResNet18 as a stable and top-performing architecture, especially relative to lightweight CNN models, while preserving competitive performance against transformer-driven approaches. Fig.9 represents The McNemar paired test results comparing ResNet18 with MobileViT-XS and ShuffleNetV2 by showcasing the distribution of discordant prediction pairs and their corresponding p-values. 

\subsection{5-Fold Cross-Validation Result}
\begin{figure}[H]
\centering
\includegraphics[width=\linewidth]{figures/cv.png}
\caption{5-fold cross-validation validation accuracy across folds (ResNet18).}
\label{fig:cv}
\end{figure}
    To further assess the consistency and generalization ability of the proposed pipeline, a 5-fold cross-validation experiment was introduced using the optimized model ResNet18. Figure X shows the gradual cross-validation accuracy of the best model against each fold. The dataset was segmented into five non-overlapping folds, where each fold was used as validation in one iteration while the remaining four folds were assigned to training. The validation accuracy across five folds stayed uniformly high, ranging roughly between 96.5 and 97.4\%, indicating stable training dynamics. The small dispersion across fold results demonstrates that the model is insensitive to data splits. This consistency validates that the learned visual features are consistent and generalizable. Overall, the CV results strengthened the reliability of ResNet18 for automated eye-disorder classification. 

\subsection{Grad-CAM Visualization}
\begin{figure}[H]
\centering
\includegraphics[width=\linewidth]{figures/grad1.png}
\label{fig:gradcam}
\end{figure}
\begin{figure}[H]
\centering
\includegraphics[width=\linewidth]{figures/grad2.png}
\caption{Grad-CAM visualizations highlighting discriminative regions.}
\label{fig:gradcam}
\end{figure}

   The GRAD-CAM visualization of the proposed model ResNet18 is shown in Figure X. Gradient-weighted Class Activation Mapping (Grad-CAM) was applied to interpret the prediction behavior of the highest-performing model, resNet18, by visualizing the spatial regions that had the strongest impact on the predictions \cite{gradcam}. As illustrated in the Grad-CAM visuals, the model repeatedly emphasized medically relevant ocular regions, especially the iris and adjacent corneal regions rather than irrelevant or background areas. For Uveitis cases, the heatmaps emphasize the central eye region, aligned with inflammation-prone regions. While Cataract samples show high attention over the lens area where lens opacity occurs. The alignment between high-activation regions and clinically recognized abnormalities indicates that the model captured meaningful disease-related features. This improves interpretability and supports the medical plausibility of the proposed system.

\subsection{Model Comparison with Previous Study}
\begin{table}[H]
\centering
\caption{Comparison with Previous Studies}
\label{tab:comparison_previous_works}
\renewcommand{\arraystretch}{1.15}
\footnotesize

\resizebox{\columnwidth}{!}{%
\begin{tabular}{|c|c|c|p{3.2cm}|}
\hline
\textbf{Study} & \textbf{Model} & \textbf{Test Acc. (\%)} & \textbf{Key Limitations} \\ \hline
Bitto et al. & ResNet50 (TL) & 95.3 & No cross-validation; no XAI \\ \hline
Zhou et al. & Vision Transformer & 95.76 & OCT-dependent; no statistical validation \\ \hline
Akça et al. & MobileViT & 99.17 & Requires OCT; lacks explainability \\ \hline
Babaqi et al. & CNN (TL) & 94.0 & Small dataset; limited robustness \\ \hline
Ayesha et al. & CNN + ViT & 94.0 & No cross-validation; weak generalization \\ \hline
\textbf{Proposed Study} & \textbf{ResNet18 (Optimized)} & \textbf{97.44} & \textbf{Cross-validation, XAI, CI, McNemar test} \\ \hline
\end{tabular}%
}
\end{table}





   
% -------------------- 5. Conclusion --------------------
\section{\textbf{Conclusion }}
This study presented a robust deep learning framework combining CNNs and transformers for multi-class eye disease classification using colored ocular images. Thorough experimental evaluation across six modern architectures revealed that ResNet18 yielded the best overall performance in terms of accuracy, consistency and generalization. The adopted preprocessing and augmentation techniques significantly improved discriminative feature quality and learning efficiency. Grad-CAM based interpretability analysis further confirmed that model predtions were guided by clinically significant eye areas. Although the current system is restricted to a limited set of eye diseases and cannot accommodate unseen or rare conditions. This confined phenomenon creates a scope in future to expand the work by including more disorder categories and multimodal clinical patient information. Overall, the framework demonstrates strong potential as a reliable and effective automated diagnosis system within its specific classes.

% -------------------- References --------------------
\bibliography{references}
\begin{thebibliography}{99}
\bibitem{1}
Bitto, A.K., Ahmed, M., Rabbi, M.R., Hossen, A., Masum, A.K.M. and Farid, D.M., 2025, July. Explainable AI Based Transfer Learning with for Multiclass Eye Disease Classification from External Eye Images. In 2025 International Conference on Quantum Photonics, Artificial Intelligence, and Networking (QPAIN) (pp. 1-7). IEEE.

\bibitem{2}
Aamir, M., Irfan, M., Ali, T., Ali, G., Shaf, A., S, A.S., Al-Beshri, A., Alasbali, T. and Mahnashi, M.H., 2020. An adoptive threshold-based multi-level deep convolutional neural network for glaucoma eye disease detection and classification. Diagnostics, 10(8), p.602.
\bibitem{3}
Zhou, Z., Niu, C., Yu, H., Zhao, J., Wang, Y. and Dai, C., 2023, March. Diagnosis of retinal diseases using the vision transformer model based on optical coherence tomography images. In SPIE-CLP Conference on Advanced Photonics 2022 (Vol. 12601, p. 1260102). SPIE
\bibitem{4}
 Akça, S., Garip, Z., Ekinci, E. and Atban, F., 2024. Automated classification of choroidal neovascularization, diabetic macular edema, and drusen from retinal OCT images using vision transformers: a comparative study. Lasers in Medical Science, 39(1), p.140.
\bibitem{5}
R.M., Bhavadharini et al. “Eye Disease Detection Using Deep Learning Models with Transfer Learning Techniques.” EAI Endorsed Trans. Scalable Inf. Syst. 12 (2024): n. pag.
\bibitem{6}
Wahab Sait, A.R., 2023. Artificial intelligence-driven eye disease classification model. Applied Sciences, 13(20), p.11437
\bibitem{7}
Mondal, S., Banerjee, S., Mukherjee, S., Ganguly, A. and Sengupta, D., 2022. Deep Classifier for Conjunctivitis–A Three-Fold Binary Approach. Int J Math Sci Comput, 2, pp.46-54.
\bibitem{8}
Nakayama, L.F., Ribeiro, L.Z., Dychiao, R.G., Zamora, Y.F., Regatieri, C.V., Celi, L.A., Silva, P., Sobrin, L. and Belfort Jr, R., 2023. Artificial intelligence in uveitis: A comprehensive review. Survey of Ophthalmology, 68(4), pp.669-677.
\bibitem{9}
Mondal, S., Banerjee, S., Mukherjee, S., Ganguly, A. and Sengupta, D., 2022. Deep Classifier for Conjunctivitis–A Three-Fold Binary Approach. Int J Math Sci Comput, 2, pp.46-54.
\bibitem{10}
Nakayama, L.F., Ribeiro, L.Z., Dychiao, R.G., Zamora, Y.F., Regatieri, C.V., Celi, L.A., Silva, P., Sobrin, L. and Belfort Jr, R., 2023. Artificial intelligence in uveitis: A comprehensive review. Survey of Ophthalmology, 68(4), pp.669-677.
\bibitem{11}
Babaqi, T., Jaradat, M., Yildirim, A.E., Al-Nimer, S.H. and Won, D., 2023. Eye disease classification using deep learning techniques. arXiv preprint arXiv:2307.10501.   
\bibitem{12}
https://ijrpr.com/uploads/V6ISSUE6/IJRPR49399.pdf
\bibitem{13}
Abdullah, A.A., Aldhahab, A. and Al Abboodi, H.M., 2024. Review of eye diseases detection and classification using deep learning techniques. In BIO Web of Conferences (Vol. 97, p. 00012). EDP Sciences. 
\bibitem{14}
Alzubaidi, L., Zhang, J., Humaidi, A.J., Al-Dujaili, A., Duan, Y., Al-Shamma, O., Santamaría, J., Fadhel, M.A., Al-Amidie, M. and Farhan, L., 2021.
\bibitem{15}
Muhammad. (2024). A Review of Transformer-Based Models for Computer Vision Tasks: Capturing Global Context and Spatial Relationships. 10.48550/arXiv.2408.15178.
\bibitem{16}
Ayesha, Noor. (2025). A Vision Transformer-Based Convolutional Neural Network for the Automated Diagnosis of Eye Diseases Using Self-Attention Mechanisms. Engineering, Technology \& Applied Science Research. 15. 24493-24497. 10.48084/etasr.10649.
\bibitem{17}
Topaloglu, Ismail. (2022). Deep Learning Based Convolutional Neural Network Structured New Image Classification Approach for Eye Disease Identification. Scientia Iranica. 10.24200/SCI.2022.58049.5537.
\bibitem{18}
Jain, Piyush. (2025). A Hybrid Deep Learning Approach for Eye Disease Detection with Integrating Adaptive CNN Layers for Improved Performance of Image Classification. Journal of Information Systems Engineering and Management. 10. 898-914. 10.52783/jisem.v10i38s.7006.
\bibitem{19}
Lin, Ankang. (2025). Efficient fusion transformer model for accurate classification of eye diseases. Scientific Reports. 15. 10.1038/s41598-025-20226-z.

\bibitem{dataset}
S. K. Das et al., “Eye Disease Image Dataset,” Mendeley Data, vol. 1, 2021. [Online]. Available: https://data.mendeley.com/datasets/n9zp473wfw/1

\bibitem{resnet}
K. He, X. Zhang, S. Ren, and J. Sun, “Deep residual learning for image recognition,” in \emph{Proc. IEEE Conf. Computer Vision and Pattern Recognition (CVPR)}, 2016, pp. 770–778.

\bibitem{vit}
A. Dosovitskiy et al., “An image is worth 16×16 words: Transformers for image recognition at scale,” in \emph{Proc. Int. Conf. Learning Representations (ICLR)}, 2021.

\bibitem{swin}
Z. Liu et al., “Swin transformer: Hierarchical vision transformer using shifted windows,” in \emph{Proc. IEEE/CVF Int. Conf. Computer Vision (ICCV)}, 2021, pp. 10012–10022.

\bibitem{mobilevit}
S. Mehta and M. Rastegari, “MobileViT: Light-weight, general-purpose, and mobile-friendly vision transformer,” in \emph{Proc. Int. Conf. Learning Representations (ICLR)}, 2022.

\bibitem{levit}
B. Graham et al., “LeViT: A vision transformer in convnet’s clothing for faster inference,” in \emph{Proc. IEEE/CVF Int. Conf. Computer Vision (ICCV)}, 2021.

\bibitem{gradcam}
R. R. Selvaraju et al., “Grad-CAM: Visual explanations from deep networks via gradient-based localization,” in \emph{Proc. IEEE Int. Conf. Computer Vision (ICCV)}, 2017, pp. 618–626.

\bibitem{clahe}
K. Zuiderveld, “Contrast limited adaptive histogram equalization,” in \emph{Graphics Gems IV}, Academic Press, 1994, pp. 474–485.

\bibitem{mcnemar}
Q. McNemar, “Note on the sampling error of the difference between correlated proportions or percentages,” \emph{Psychometrika}, vol. 12, no. 2, pp. 153–157, 1947.

\bibitem{medical_metrics}
T. Fawcett, “An introduction to ROC analysis,” \emph{Pattern Recognition Letters}, vol. 27, no. 8, pp. 861–874, 2006.


\end{thebibliography}

\end{document}
